% =============================================================================
% Section 3: Main theorems
% Drafted FIRST per user directive: "lock the theorem stack before polishing prose"
% =============================================================================

\section{Main results}
\label{sec:main}

The contribution of this paper consists of six theorems whose
collective effect is to characterize the structure of the residual
frontier of Erd\H{o}s Problem \#647 and to certify that the natural
finite-local closure paradigm is exhausted at the parameter regimes
tested. We state the theorems precisely here; subsequent sections
provide proofs, computational certificates, and contextual analysis.

\subsection{Stage-1 sieve reduction}
\label{sec:thm:stage1}

\begin{theorem}[Stage-1 sieve reduction; Lean-certified]
\label{thm:stage1-reduction}
Let $n > 24$ be an integer satisfying
\[
\max_{1 \le m < n} \bigl( m + \tauf(m) \bigr) \le n + 2 .
\]
Then $n \equiv r \pmod{M}$, where $M = 46189 = 11 \cdot 13 \cdot 17 \cdot 19$
and $r \in \Ropen$ for an explicit set $\Ropen \subset \Z / M\Z$ with
$|\Ropen| = 41$. The set $\Ropen$ is fully Lean-formalized
(Mathlib import: \texttt{Erdos647ResiduePartitionStage1}); a
machine-verified proof of the reduction
is provided in \cref{sec:structural-reduction}.
\end{theorem}

\subsection{Universal $v_{23}$ partial closure}
\label{sec:thm:v23}

\begin{theorem}[Universal $v_{23}$ partial closure; Lean-certified]
\label{thm:v23-partial-closure}
For each of the $96$ residue classes $r$ modulo $M$ surviving the
Stage-1 sieve, there exists an explicit
sub-AP $r' \pmod{M \cdot d}$ (for a fixed $d = 1058 = 2 \cdot 23^2$)
such that no $n \equiv r' \pmod{M d}$ satisfies the
$\tauf$-barrier inequality. The map $r \mapsto r'$ is computer-checked
and Lean-formalized
(\texttt{lean/generated/V23All96.lean}, $96$ instances of the
\texttt{hensel\_rootline\_closure} primitive).

This partial closure is universal in the sense that it eliminates a
common $v_{23}$-rich sublattice across all $96$ classes, and is
independent of the further reductions in
\cref{thm:stage1-reduction,thm:143-collapse}.
\end{theorem}

\begin{remark}
The $v_{23}$ closure does not by itself reduce the open count below
$96$, but it removes a structurally common sub-AP that would otherwise
need to be re-handled at every later layer. We retain the full
$\Ropen$ set in subsequent statements; the $v_{23}$ partial closure
is a separately citable Lean theorem that constrains the
$v_{23}$-richness of any potential counterexample.
\end{remark}

\subsection{The 143-collapse}
\label{sec:thm:143-collapse}

\begin{theorem}[143-collapse; Lean-certified]
\label{thm:143-collapse}
Every $r \in \Ropen$ satisfies $11 \mid r$ and $13 \mid r$, hence
$143 \mid r$. Equivalently, for every potential counterexample $n$
to Erd\H{o}s~\#647 satisfying $n > 24$, there exists a unique
decomposition
\[
n = 143 \cdot Y, \qquad Y \in \Z, \qquad Y \bmod 323 \in \Uset ,
\]
where $\Uset \subset \Z / 323\Z$ is the explicit 41-element set
$\{0, 6, 9, 12, 17, 18, 34, 37, 43, 57, 63, 69, 74, 85, 88, 91, 94,
114, 119, 120, 131, 139, 142, 145, 148, 170, 171, 176, 190, 196, 199,
205, 207, 221, 224, 227, 247, 256, 264, 272, 309\}$.

Lean formalization: \texttt{Erdos647\_143Collapse.lean}
(theorems \texttt{openResidues41\_143\_collapse} and
\texttt{U\_41\_explicit}; both proved by \texttt{native\_decide}).
\end{theorem}

\begin{remark}
\Cref{thm:143-collapse} replaces the original modulus $M = 46189$
with the much cleaner normalization $q = 323 = 17 \cdot 19$. We will
use the equivalent rewriting
\[
2520 \cdot n - k \;=\; 360360 \cdot Y - k, \qquad
360360 = \lcm(1, 2, \ldots, 15),
\]
throughout, particularly in the analysis of forced low-prime divisors
of small shifts of any potential counterexample (\cref{sec:csp-audit}).
\end{remark}

\begin{remark}[Notation: search variable vs.\ candidate]
\label{rem:n-convention}
Throughout \cref{thm:stage1-reduction,thm:143-collapse,sec:thm:v23} the
symbol $n$ denotes the \emph{search variable} entering the Stage-1 sieve
analysis, equal to $n_{\rm act} / 2520$ where $n_{\rm act}$ is the
hypothetical \#647 counterexample of \cref{prob:e647}; this matches the
Lean signature
\texttt{bridge\_isErdos647\_to\_sieve(N)} taking the candidate as
$2520 \cdot N$ (\texttt{lean/Bridge2.lean}, line~1013). The residue
restriction ``$n \equiv r \pmod M$, $r \in \Ropen$'' should accordingly
be read as ``$n_{\rm act} / 2520 \equiv r \pmod M$.'' The full-divisibility
corollary stated in \cref{cor:full-divisibility} below packages this
explicitly for downstream citation.
\end{remark}

\subsection{Finite-local insufficiency theorem}
\label{sec:thm:ple-barrier}

\begin{definition}[Pointwise-local-existential class]
\label{def:ple}
The class \PLEclass\ of \emph{pointwise-local-existential
forced-divisor methods} consists of all closure primitives that decide
the existence of a counterexample $n \equiv r \pmod M$ ($r \in \Ropen$)
based on local data at a finite set of primes, with no averaged,
density, or sieve-weight component. Concretely, \PLEclass\ includes:

\begin{enumerate}[label=(\roman*), leftmargin=*]
\item Cycle-3 valuation leaves;
\item Multilevel DivisorMask primitives (free-prime, joint $p$-adic,
  exact-valuation $\mu$-DivisorMask);
\item Hensel-balanced prime-residual kernels (balanced and unbalanced)
  with the per-shift admissibility condition $2\tauf(D_k) \le k+2$
  for $k \le K_0$;
\item Fixed-kernel interval structural-kill detectors;
\item Kernel-space CSPs varying the kernel base $a$ subject to the
  Hensel-admissibility condition;
\item Full safe-set CSPs combining kernel-base variables with
  auxiliary safe-set residues for primes $\ell \le B$.
\end{enumerate}

A \emph{closure certificate} in this class is a finite-local proof
that no $n \equiv r \pmod M$ satisfies the $\tauf$-barrier inequality
beyond an explicit threshold $n_0$.
\end{definition}

\begin{theorem}[Finite-local insufficiency; SAT-certified]
\label{thm:ple-insufficiency}
For every sentinel residue $r \in \{1716, 4862\}$ and every parameter
choice $(K_0, B, K_{\rm scan})$ tested with $K_0 \in \{95, 200, 1021\}$,
$B \le 1000$, $K_{\rm scan} \le 10^5$, the encoded \PLEclass\ instance
$\Phi_r(K_0, B, K_{\rm scan})$ is satisfiable. Hence no \PLEclass\
primitive in the encoded freedom yields a closure certificate at the
sentinel residues.

The satisfiability of $\Phi_r$ is certified by a CaDiCaL-generated
LRAT proof, verified by \texttt{cake\_lpr} (CakeML-based LRAT
checker), and importable into Lean via
\texttt{Std.Tactic.BVDecide}.
\end{theorem}

\begin{remark}[Quantifier honesty]
\label{rem:per-kernel-class}
At parameter $K_{\rm scan} \le 10^7$, the structural-kill detector
identifies four residues with per-kernel-class certificates:
$r \in \{858, 1716, 2431, 4862\}$. Per the quantifier audit
(\cref{sec:csp-audit}, summarized from
\texttt{megak\_closure\_quantifier\_audit\_20260425.md}),
\emph{none of these certificates extend to a base-AP-universal
closure}. The mega-shift artifacts at $k = 2520 r$ for $r \in \{858, 2431\}$
arise from the structural identity $F_k(s) = A \cdot s$ when $c = 0$,
and apply only to the JSON kernel CRT class chosen by the solver,
not to all $n$ in the residue class. Consequently \emph{no residue is
unconditionally closed by \PLEclass\ at the tested parameters}.
\end{remark}

\subsection{Conditional barrier theorem}
\label{sec:thm:conditional-barrier}

\begin{theorem}[Conditional barrier; Schinzel-H form]
\label{thm:conditional-barrier}
Assume Schinzel's Hypothesis H. For every $u \in \Uset$ and every
fixed $K \ge 1$, the polynomial system $\Sigma_u(K)$ encoding the
\Cref{thm:stage1-reduction} normalization restricted to the
$\tauf$-barrier window of length $K$ is admissible (no fixed prime
divisor) and the Bateman--Horn singular series
$\mathfrak{S}(\Sigma_u, K)$ is positive. Conjecturally, each of the
41 residue classes contains infinitely many $n$ satisfying the
$K$-window $\tauf$-barrier inequality, with predicted density
\[
N_u(X; K) \;\sim\; \mathfrak{S}(\Sigma_u, K) \cdot \frac{X}{(\log X)^{\deg \Sigma_u}}
\quad \text{as } X \to \infty.
\]
\end{theorem}

\begin{remark}
We do \emph{not} claim this implies closure or non-closure of
\#647 itself; \cref{thm:conditional-barrier} establishes the
\emph{infinite-survivor} structural picture conditional on H, in
explicit contradistinction to a \PLEclass\ closure certificate. The
combination of \cref{thm:ple-insufficiency,thm:conditional-barrier}
constitutes the \emph{barrier} of the title: under H, no \PLEclass\
finite-local forced-divisor closure of the residual frontier exists,
because the survivor set is genuinely infinite at every $K$.
\end{remark}

\begin{remark}[Quantitative empirical anchor]
\label{rem:conditional-empirical-anchor}
The empirical multipoint $\tauf$-correlation analysis of
\cref{sec:empirical-correlations} provides quantitative content to
the singular-series prediction underlying
\cref{thm:conditional-barrier}. The pairwise correlation matrix
across $k \in [1, 100]$ shifts shows a four-class lag-correlation
decomposition (\cref{tab:correlation-decomposition}) governed by
$\gcd(d, M^*)$ with $M^* = 2520\, M$, and the per-class marginal
$\bar\tauf \approx \log n$ (within $3\%$ across all $41$ classes)
matches the Dirichlet asymptotic. The empirical pattern is consistent
with, and gives quantitative empirical content to, the
Schinzel-conditional characterization of surviving classes: the
framework reduction preserves the Bateman--Horn singular series
locally and reshapes the lag-correlation structure in a
sign-dependent way governed by $v_2$-parity.
\end{remark}

\begin{remark}[On the GRH-46189 + Hensel + Schinzel-H formulation]
\label{rem:no-cgk-closure}
A natural conjecture is that adding GRH at $q = 46189$ to the
Hensel admissibility lemma + Schinzel-H promotes the conditional
barrier to a conditional closure. We do \emph{not} make this claim:
per
\cref{sec:obstructions} (and the synthesis of the contemporary
correlation-of-arithmetic-functions literature in
\texttt{closure\_via\_correlation\_methods\_synthesis\_20260425.md}),
this would additionally require a pointwise version of the
Conrey--Gonek--Keating asymptotic at $q = 46189$, which is open
(open even unconditionally for $k = l = 3$ in the divisor
correlation $\sum d_3(n) d_3(n+h)$; see Matomäki--Radziwiłł--Tao
2019).
\end{remark}

\subsection{Empirical verification}
\label{sec:thm:empirical}

\begin{proposition}[Empirical floor at the principal residue]
\label{prop:empirical-r0}
No counterexample to Erd\H{o}s~\#647 exists at the principal
residue $r = 0$ contiguously for $n \le 3.34 \times 10^{21}$.
Additional non-contiguous islands of exhaustive verification extend
to $n \approx 1.34 \times 10^{22}$ at $r = 0$. Two independent
$13$-fold simultaneous-shape-prime witnesses at
$t \in \{1.17 \times 10^{13}, 1.15 \times 10^{14}\}$
(corresponding to $n \approx 10^{21}$ and $n \approx 10^{22}$
respectively) are confirmed and ruled out as counterexamples by the
non-shape shift $k = 7$ in both cases.

The 40 non-principal residue classes in $\Ropen$ are not yet
empirically verified at scale; this is identified as Open
Problem~\ref{prob:non-principal-empirical}.
\end{proposition}

\begin{remark}[Honest scope on the empirical bound]
The bound $n \le 3.34 \times 10^{21}$ is the contiguous exhaustive
floor at $r = 0$ only. We do \emph{not} claim
``no counterexample below $10^{22}$'' --- the gaps in the witness
search above $3.34 \times 10^{21}$ remain (see
\cref{sec:empirical-detail}), and the non-principal classes are
unverified empirically. The two $13/13$ witnesses are hard-case
samples at the predicted Bateman--Horn density, not exhaustive
verification.
\end{remark}

% =============================================================================
% End of section 3
% =============================================================================
